\topic{从电气条件到指令提交}

硬件书从电压、电流和时序进入逻辑，本书必须保留同一条下落路径。电源管理器
先让各路供电进入容差窗口，参考时钟经过 PLL 倍频并稳定，复位树再按时钟域
同步释放。处理器前端取得复位地址，平台译码选择 ROM，存储器返回若干字节，
取指队列识别指令边界，译码器生成内部操作，执行单元改变架构状态，最后退休
逻辑使变化对软件可见。

\begin{osgridlongtable}{L{0.18\textwidth}L{0.28\textwidth}L{0.40\textwidth}}
\textbf{层次} & \textbf{建立的事实} & \textbf{失效时的外部表现}\\ \hline
\endfirsthead
\textbf{层次} & \textbf{建立的事实} & \textbf{失效时的外部表现}\\ \hline
\endhead
供电 & 电压、纹波和瞬态电流满足器件范围 & 无法复位、随机重启或边沿误判\\ \hline
时钟 & 周期和抖动满足采样窗口 & 状态机停滞或跨域采样不稳定\\ \hline
复位 & 各时钟域从合法状态开始 & 控制逻辑进入未定义组合\\ \hline
地址译码 & 唯一存储器响应复位地址 & 读到全一、总线错误或错误设备\\ \hline
取指与译码 & 字节被分割为合法 x86 指令 & 非法指令或错误控制流\\ \hline
执行与退休 & 指令效果按 ISA 顺序可观察 & 异常、停机或架构状态不一致\\ \hline
\end{osgridlongtable}

软件文档常从复位向量开始，是因为电气层已经由平台规范和器件实现承担；这不
表示前置过程不存在。本项目使用 QEMU TCG 后，真实电压、PLL 锁定和板级信号
被模型化为确定的初始状态。模型边界让实验可重复，也限定了证据范围：QEMU
启动成功证明软件满足所选机器模型，不证明任意真实主板的电源完整性。

\topic{同步数字系统依赖建立时间与保持时间}
组合逻辑需要传播时间，触发器在时钟有效沿附近需要输入稳定。建立时间要求数据
在采样沿之前到达，保持时间要求数据在采样沿之后继续稳定。一个同步路径满足：
\[
  T_{\mathrm{clock}} \ge
  T_{\mathrm{clk\rightarrow q}}+
  T_{\mathrm{comb,max}}+
  T_{\mathrm{setup}}+
  T_{\mathrm{skew}}.
\]
如果目标频率提高，允许的组合传播窗口缩短，设计需要流水化或更快电路。操作
系统看不到每个门延迟，却依赖硬件在时钟约束内兑现 ISA。超频、供电不足和温度
变化之所以能造成“软件随机崩溃”，正是物理时序余量被侵蚀。

跨时钟域信号还可能在采样边沿附近变化，使触发器进入亚稳态。硬件用同步器、
握手或异步 FIFO 降低传播到系统状态的概率。软件面对的 UART、磁盘和中断
控制器本质上都是独立时钟域；状态位和队列就是硬件向软件暴露的同步边界。

\topic{处理器执行与存储层次}

ISA 可以用“取指、译码、执行”解释语义，但现代处理器会预取、乱序执行和推测
分支。为了保持软件模型，退休阶段按程序顺序提交结果，异常则表现为精确边界：
异常指令之前的效果已经提交，之后的效果对架构不可见。

启动代码依赖的是这个架构承诺，不依赖流水线级数。即使 QEMU TCG 没有模拟
真实分支预测器，它也必须给出相同的寄存器、内存、I/O 和异常结果。区分架构
与微架构可以避免两种错误：把某代 CPU 的性能现象写成正确性要求，或者把内存
顺序、缓存一致性等真正的架构约束误认为实现细节。

\topic{同一份字节经过缓存以后由谁负责写回}
寄存器通常在核心内部，缓存保存最近使用的内存副本，DRAM 保存大容量易失
数据，Flash 或 ROM 保存掉电后仍存在的启动字节。越接近执行单元的层次容量
越小、延迟越低。缓存行常以几十字节为单位搬运，而软件地址仍按字节编号。

\begin{osgridlongtable}{L{0.18\textwidth}L{0.19\textwidth}L{0.25\textwidth}L{0.27\textwidth}}
\textbf{层次} & \textbf{典型粒度} & \textbf{主要管理者} & \textbf{启动阶段作用}\\ \hline
\endfirsthead
\textbf{层次} & \textbf{典型粒度} & \textbf{主要管理者} & \textbf{启动阶段作用}\\ \hline
\endhead
寄存器 & 8--64 位 & 指令与 ABI & 保存当前控制和数据状态\\ \hline
缓存 & 缓存行 & 硬件一致性协议 & 加速取指与 RAM 访问\\ \hline
DRAM & 字节寻址、突发传输 & 内存控制器与内核 & 栈、页表、内核和进程\\ \hline
ROM/Flash & 镜像与擦写块 & 平台与固件构建 & 提供复位后立即可读的代码\\ \hline
磁盘 & 扇区与块 & 控制器、驱动和文件系统 & 保存后续加载阶段与持久数据\\ \hline
\end{osgridlongtable}

缓存不会创造新的数据所有权。内核仍要记录物理页属于谁，文件系统仍要记录
磁盘块属于哪个对象。存储层次只改变副本位置与完成时刻，因此后续必须讨论
缓存一致性、写回、屏障和持久化语义。

\topic{固件、引导器与内核的责任分层}

早期机器由操作员直接装入程序，后来 ROM 监控程序提供基本装载和诊断。PC
BIOS 把硬件初始化与若干通用服务固化为接口，磁盘引导扇区再装入更复杂系统。
UEFI 进一步使用可扩展驱动、文件协议和统一可执行格式。GRUB、Limine 等现代
引导器可以替内核处理大量平台细节。

这些成熟方案适合生产系统，却会隐藏学习目标。本项目不否定它们，而是把责任
重新取回：自研 ROM 建立最小环境，自研装载阶段读取磁盘并切换模式，自研
内核接管异常、内存和设备。亲自实现以后，才有能力准确评价现成固件与引导器
提供了哪些保证。

\section{操作系统抽象、状态与权限}

进程不是一个编号，文件也不是一个路径字符串。进程抽象由页表、线程现场、
打开对象、权限和生命周期状态共同维持；文件抽象由 inode、目录项、缓存、
块映射和持久化规则共同维持。任何抽象都必须对应可检查的不变量。

保护是操作系统区别于普通运行库的关键。处理器特权级限制敏感指令，页表限制
地址访问，系统调用验证不可信参数，内核对象记录所有权。若某一层只提供方便
接口而没有权限边界，它仍可能是运行时或库，却不是完整的操作系统隔离。

\topic{运行第一个实验时能观察到什么}
当前工程能够直接证明的最早硬件链包括复位地址、ROM 字节、入口指令和 UART
端口事务。电源与时钟部分属于理解真实机器所需的背景，在 QEMU 中没有模拟
电压波形。书中因此把“规范依据”“模型行为”“当前测试证据”分开表述。

工程证据从最终产物向上闭合：ROM 固定偏移包含复位跳转，反汇编显示目标入口，
QEMU 串口留下阶段标记，失败镜像证明轮询能终止，CI 在干净宿主重新执行。
只有这条链上的事实才被当前版本声明为已实现。
